% ============================================================
%  Shared preamble for all course notes
%  Loaded from each course's main.tex via % ============================================================
%  Shared preamble for all course notes
%  Loaded from each course's main.tex via % ============================================================
%  Shared preamble for all course notes
%  Loaded from each course's main.tex via % ============================================================
%  Shared preamble for all course notes
%  Loaded from each course's main.tex via \input{../common/preamble}
% ============================================================

% ---------- Page layout & fonts ----------
\usepackage[utf8]{inputenc}
\usepackage[T1]{fontenc}
\usepackage{lmodern}
\usepackage{microtype}
\usepackage[a4paper, margin=2.5cm]{geometry}
\usepackage{parskip}               % space between paragraphs, no indent
\usepackage{fancyhdr}

% ---------- Math ----------
\usepackage{amsmath, amssymb, amsthm}
\usepackage{mathtools}
\usepackage{bm}

% ---------- Figures, tables, plots ----------
\usepackage{graphicx}
\usepackage{float}
\usepackage{caption}
\usepackage{subcaption}
\usepackage{booktabs}
\usepackage{tikz}
\usetikzlibrary{arrows.meta, positioning, shapes, calc, fit}
\usepackage{pgfplots}
\pgfplotsset{compat=1.17}
\graphicspath{{figures/}}

% ---------- Lists, code, colored boxes ----------
\usepackage{enumitem}
\usepackage{xcolor}
\usepackage{listings}
\usepackage[most]{tcolorbox}

% ---------- References ----------
\usepackage[numbers]{natbib}
\usepackage[colorlinks=true, linkcolor=blue!60!black,
            citecolor=green!50!black, urlcolor=blue!70!black]{hyperref}
\usepackage[capitalise, noabbrev]{cleveref}   % must come after hyperref

% ---------- Header / footer ----------
\pagestyle{fancy}
\fancyhf{}
\fancyhead[L]{\nouppercase{\leftmark}}
\fancyhead[R]{\coursetitle}
\fancyfoot[C]{\thepage}
\setlength{\headheight}{14.5pt}

% ---------- Theorem-like environments ----------
\theoremstyle{definition}
\newtheorem{definition}{Definition}[chapter]
\newtheorem{example}{Example}[chapter]
\newtheorem{exercise}{Exercise}[chapter]

\theoremstyle{plain}
\newtheorem{theorem}{Theorem}[chapter]
\newtheorem{lemma}[theorem]{Lemma}
\newtheorem{proposition}[theorem]{Proposition}
\newtheorem{corollary}[theorem]{Corollary}

\theoremstyle{remark}
\newtheorem*{remark}{Remark}
\newtheorem*{note}{Note}

% ---------- Highlight boxes ----------
% \begin{keyidea} ... \end{keyidea}   -- main takeaway of a lecture
% \begin{intuition} ... \end{intuition}
\newtcolorbox{keyidea}[1][Key Idea]{
  colback=blue!4, colframe=blue!50!black, fonttitle=\bfseries,
  title=#1, breakable}
\newtcolorbox{intuition}[1][Intuition]{
  colback=green!4, colframe=green!40!black, fonttitle=\bfseries,
  title=#1, breakable}
\newtcolorbox{warning}[1][Caution]{
  colback=red!4, colframe=red!50!black, fonttitle=\bfseries,
  title=#1, breakable}

% ---------- Code listings (Python by default) ----------
\lstdefinestyle{notes}{
  language=Python,
  basicstyle=\ttfamily\small,
  keywordstyle=\color{blue!70!black}\bfseries,
  commentstyle=\color{green!45!black}\itshape,
  stringstyle=\color{red!60!black},
  numbers=left, numberstyle=\tiny\color{gray},
  frame=single, rulecolor=\color{gray!40},
  backgroundcolor=\color{gray!5},
  breaklines=true, showstringspaces=false, tabsize=4,
  captionpos=b
}
\lstset{style=notes}

% ---------- Shared math notation ----------
\input{../common/macros}

% ============================================================

% ---------- Page layout & fonts ----------
\usepackage[utf8]{inputenc}
\usepackage[T1]{fontenc}
\usepackage{lmodern}
\usepackage{microtype}
\usepackage[a4paper, margin=2.5cm]{geometry}
\usepackage{parskip}               % space between paragraphs, no indent
\usepackage{fancyhdr}

% ---------- Math ----------
\usepackage{amsmath, amssymb, amsthm}
\usepackage{mathtools}
\usepackage{bm}

% ---------- Figures, tables, plots ----------
\usepackage{graphicx}
\usepackage{float}
\usepackage{caption}
\usepackage{subcaption}
\usepackage{booktabs}
\usepackage{tikz}
\usetikzlibrary{arrows.meta, positioning, shapes, calc, fit}
\usepackage{pgfplots}
\pgfplotsset{compat=1.17}
\graphicspath{{figures/}}

% ---------- Lists, code, colored boxes ----------
\usepackage{enumitem}
\usepackage{xcolor}
\usepackage{listings}
\usepackage[most]{tcolorbox}

% ---------- References ----------
\usepackage[numbers]{natbib}
\usepackage[colorlinks=true, linkcolor=blue!60!black,
            citecolor=green!50!black, urlcolor=blue!70!black]{hyperref}
\usepackage[capitalise, noabbrev]{cleveref}   % must come after hyperref

% ---------- Header / footer ----------
\pagestyle{fancy}
\fancyhf{}
\fancyhead[L]{\nouppercase{\leftmark}}
\fancyhead[R]{\coursetitle}
\fancyfoot[C]{\thepage}
\setlength{\headheight}{14.5pt}

% ---------- Theorem-like environments ----------
\theoremstyle{definition}
\newtheorem{definition}{Definition}[chapter]
\newtheorem{example}{Example}[chapter]
\newtheorem{exercise}{Exercise}[chapter]

\theoremstyle{plain}
\newtheorem{theorem}{Theorem}[chapter]
\newtheorem{lemma}[theorem]{Lemma}
\newtheorem{proposition}[theorem]{Proposition}
\newtheorem{corollary}[theorem]{Corollary}

\theoremstyle{remark}
\newtheorem*{remark}{Remark}
\newtheorem*{note}{Note}

% ---------- Highlight boxes ----------
% \begin{keyidea} ... \end{keyidea}   -- main takeaway of a lecture
% \begin{intuition} ... \end{intuition}
\newtcolorbox{keyidea}[1][Key Idea]{
  colback=blue!4, colframe=blue!50!black, fonttitle=\bfseries,
  title=#1, breakable}
\newtcolorbox{intuition}[1][Intuition]{
  colback=green!4, colframe=green!40!black, fonttitle=\bfseries,
  title=#1, breakable}
\newtcolorbox{warning}[1][Caution]{
  colback=red!4, colframe=red!50!black, fonttitle=\bfseries,
  title=#1, breakable}

% ---------- Code listings (Python by default) ----------
\lstdefinestyle{notes}{
  language=Python,
  basicstyle=\ttfamily\small,
  keywordstyle=\color{blue!70!black}\bfseries,
  commentstyle=\color{green!45!black}\itshape,
  stringstyle=\color{red!60!black},
  numbers=left, numberstyle=\tiny\color{gray},
  frame=single, rulecolor=\color{gray!40},
  backgroundcolor=\color{gray!5},
  breaklines=true, showstringspaces=false, tabsize=4,
  captionpos=b
}
\lstset{style=notes}

% ---------- Shared math notation ----------
% ============================================================
%  Shared math notation for ML / DL notes
% ============================================================

% ---------- Number sets ----------
\newcommand{\R}{\mathbb{R}}
\newcommand{\N}{\mathbb{N}}
\newcommand{\Z}{\mathbb{Z}}

% ---------- Probability ----------
\newcommand{\E}{\mathbb{E}}
\newcommand{\Prob}{\mathbb{P}}
\DeclareMathOperator{\Var}{Var}
\DeclareMathOperator{\Cov}{Cov}
\newcommand{\KL}[2]{D_{\mathrm{KL}}\!\left(#1 \,\|\, #2\right)}
\newcommand{\Normal}{\mathcal{N}}
\newcommand{\iid}{\overset{\text{i.i.d.}}{\sim}}

% ---------- Linear algebra ----------
\renewcommand{\vec}[1]{\bm{#1}}          % vectors: \vec{x}
\newcommand{\mat}[1]{\bm{#1}}            % matrices: \mat{W}
\newcommand{\T}{^{\top}}                 % transpose: \vec{x}\T
\DeclareMathOperator{\tr}{tr}
\DeclareMathOperator{\rank}{rank}
\DeclareMathOperator{\diag}{diag}
\DeclarePairedDelimiter{\norm}{\lVert}{\rVert}
\DeclarePairedDelimiter{\abs}{\lvert}{\rvert}
\DeclarePairedDelimiter{\inner}{\langle}{\rangle}

% ---------- Optimization / calculus ----------
\DeclareMathOperator*{\argmin}{arg\,min}
\DeclareMathOperator*{\argmax}{arg\,max}
\newcommand{\grad}{\nabla}
\newcommand{\pd}[2]{\frac{\partial #1}{\partial #2}}

% ---------- ML-specific ----------
\newcommand{\dataset}{\mathcal{D}}
\newcommand{\loss}{\mathcal{L}}
\newcommand{\hypo}{\mathcal{H}}
\newcommand{\params}{\bm{\theta}}
\newcommand{\X}{\mathcal{X}}
\newcommand{\Y}{\mathcal{Y}}
\newcommand{\indicator}[1]{\mathbb{1}\!\left[#1\right]}
\DeclareMathOperator{\sigmoid}{\sigma}
\DeclareMathOperator{\softmax}{softmax}
\DeclareMathOperator{\ReLU}{ReLU}
\DeclareMathOperator{\sign}{sign}


% ============================================================

% ---------- Page layout & fonts ----------
\usepackage[utf8]{inputenc}
\usepackage[T1]{fontenc}
\usepackage{lmodern}
\usepackage{microtype}
\usepackage[a4paper, margin=2.5cm]{geometry}
\usepackage{parskip}               % space between paragraphs, no indent
\usepackage{fancyhdr}

% ---------- Math ----------
\usepackage{amsmath, amssymb, amsthm}
\usepackage{mathtools}
\usepackage{bm}

% ---------- Figures, tables, plots ----------
\usepackage{graphicx}
\usepackage{float}
\usepackage{caption}
\usepackage{subcaption}
\usepackage{booktabs}
\usepackage{tikz}
\usetikzlibrary{arrows.meta, positioning, shapes, calc, fit}
\usepackage{pgfplots}
\pgfplotsset{compat=1.17}
\graphicspath{{figures/}}

% ---------- Lists, code, colored boxes ----------
\usepackage{enumitem}
\usepackage{xcolor}
\usepackage{listings}
\usepackage[most]{tcolorbox}

% ---------- References ----------
\usepackage[numbers]{natbib}
\usepackage[colorlinks=true, linkcolor=blue!60!black,
            citecolor=green!50!black, urlcolor=blue!70!black]{hyperref}
\usepackage[capitalise, noabbrev]{cleveref}   % must come after hyperref

% ---------- Header / footer ----------
\pagestyle{fancy}
\fancyhf{}
\fancyhead[L]{\nouppercase{\leftmark}}
\fancyhead[R]{\coursetitle}
\fancyfoot[C]{\thepage}
\setlength{\headheight}{14.5pt}

% ---------- Theorem-like environments ----------
\theoremstyle{definition}
\newtheorem{definition}{Definition}[chapter]
\newtheorem{example}{Example}[chapter]
\newtheorem{exercise}{Exercise}[chapter]

\theoremstyle{plain}
\newtheorem{theorem}{Theorem}[chapter]
\newtheorem{lemma}[theorem]{Lemma}
\newtheorem{proposition}[theorem]{Proposition}
\newtheorem{corollary}[theorem]{Corollary}

\theoremstyle{remark}
\newtheorem*{remark}{Remark}
\newtheorem*{note}{Note}

% ---------- Highlight boxes ----------
% \begin{keyidea} ... \end{keyidea}   -- main takeaway of a lecture
% \begin{intuition} ... \end{intuition}
\newtcolorbox{keyidea}[1][Key Idea]{
  colback=blue!4, colframe=blue!50!black, fonttitle=\bfseries,
  title=#1, breakable}
\newtcolorbox{intuition}[1][Intuition]{
  colback=green!4, colframe=green!40!black, fonttitle=\bfseries,
  title=#1, breakable}
\newtcolorbox{warning}[1][Caution]{
  colback=red!4, colframe=red!50!black, fonttitle=\bfseries,
  title=#1, breakable}

% ---------- Code listings (Python by default) ----------
\lstdefinestyle{notes}{
  language=Python,
  basicstyle=\ttfamily\small,
  keywordstyle=\color{blue!70!black}\bfseries,
  commentstyle=\color{green!45!black}\itshape,
  stringstyle=\color{red!60!black},
  numbers=left, numberstyle=\tiny\color{gray},
  frame=single, rulecolor=\color{gray!40},
  backgroundcolor=\color{gray!5},
  breaklines=true, showstringspaces=false, tabsize=4,
  captionpos=b
}
\lstset{style=notes}

% ---------- Shared math notation ----------
% ============================================================
%  Shared math notation for ML / DL notes
% ============================================================

% ---------- Number sets ----------
\newcommand{\R}{\mathbb{R}}
\newcommand{\N}{\mathbb{N}}
\newcommand{\Z}{\mathbb{Z}}

% ---------- Probability ----------
\newcommand{\E}{\mathbb{E}}
\newcommand{\Prob}{\mathbb{P}}
\DeclareMathOperator{\Var}{Var}
\DeclareMathOperator{\Cov}{Cov}
\newcommand{\KL}[2]{D_{\mathrm{KL}}\!\left(#1 \,\|\, #2\right)}
\newcommand{\Normal}{\mathcal{N}}
\newcommand{\iid}{\overset{\text{i.i.d.}}{\sim}}

% ---------- Linear algebra ----------
\renewcommand{\vec}[1]{\bm{#1}}          % vectors: \vec{x}
\newcommand{\mat}[1]{\bm{#1}}            % matrices: \mat{W}
\newcommand{\T}{^{\top}}                 % transpose: \vec{x}\T
\DeclareMathOperator{\tr}{tr}
\DeclareMathOperator{\rank}{rank}
\DeclareMathOperator{\diag}{diag}
\DeclarePairedDelimiter{\norm}{\lVert}{\rVert}
\DeclarePairedDelimiter{\abs}{\lvert}{\rvert}
\DeclarePairedDelimiter{\inner}{\langle}{\rangle}

% ---------- Optimization / calculus ----------
\DeclareMathOperator*{\argmin}{arg\,min}
\DeclareMathOperator*{\argmax}{arg\,max}
\newcommand{\grad}{\nabla}
\newcommand{\pd}[2]{\frac{\partial #1}{\partial #2}}

% ---------- ML-specific ----------
\newcommand{\dataset}{\mathcal{D}}
\newcommand{\loss}{\mathcal{L}}
\newcommand{\hypo}{\mathcal{H}}
\newcommand{\params}{\bm{\theta}}
\newcommand{\X}{\mathcal{X}}
\newcommand{\Y}{\mathcal{Y}}
\newcommand{\indicator}[1]{\mathbb{1}\!\left[#1\right]}
\DeclareMathOperator{\sigmoid}{\sigma}
\DeclareMathOperator{\softmax}{softmax}
\DeclareMathOperator{\ReLU}{ReLU}
\DeclareMathOperator{\sign}{sign}


% ============================================================

% ---------- Page layout & fonts ----------
\usepackage[utf8]{inputenc}
\usepackage[T1]{fontenc}
\usepackage{lmodern}
\usepackage{microtype}
\usepackage[a4paper, margin=2.5cm]{geometry}
\usepackage{parskip}               % space between paragraphs, no indent
\usepackage{fancyhdr}

% ---------- Math ----------
\usepackage{amsmath, amssymb, amsthm}
\usepackage{mathtools}
\usepackage{bm}

% ---------- Figures, tables, plots ----------
\usepackage{graphicx}
\usepackage{float}
\usepackage{caption}
\usepackage{subcaption}
\usepackage{booktabs}
\usepackage{tikz}
\usetikzlibrary{arrows.meta, positioning, shapes, calc, fit}
\usepackage{pgfplots}
\pgfplotsset{compat=1.17}
\graphicspath{{figures/}}

% ---------- Lists, code, colored boxes ----------
\usepackage{enumitem}
\usepackage{xcolor}
\usepackage{listings}
\usepackage[most]{tcolorbox}

% ---------- References ----------
\usepackage[numbers]{natbib}
\usepackage[colorlinks=true, linkcolor=blue!60!black,
            citecolor=green!50!black, urlcolor=blue!70!black]{hyperref}
\usepackage[capitalise, noabbrev]{cleveref}   % must come after hyperref

% ---------- Header / footer ----------
\pagestyle{fancy}
\fancyhf{}
\fancyhead[L]{\nouppercase{\leftmark}}
\fancyhead[R]{\coursetitle}
\fancyfoot[C]{\thepage}
\setlength{\headheight}{14.5pt}

% ---------- Theorem-like environments ----------
\theoremstyle{definition}
\newtheorem{definition}{Definition}[chapter]
\newtheorem{example}{Example}[chapter]
\newtheorem{exercise}{Exercise}[chapter]

\theoremstyle{plain}
\newtheorem{theorem}{Theorem}[chapter]
\newtheorem{lemma}[theorem]{Lemma}
\newtheorem{proposition}[theorem]{Proposition}
\newtheorem{corollary}[theorem]{Corollary}

\theoremstyle{remark}
\newtheorem*{remark}{Remark}
\newtheorem*{note}{Note}

% ---------- Highlight boxes ----------
% \begin{keyidea} ... \end{keyidea}   -- main takeaway of a lecture
% \begin{intuition} ... \end{intuition}
\newtcolorbox{keyidea}[1][Key Idea]{
  colback=blue!4, colframe=blue!50!black, fonttitle=\bfseries,
  title=#1, breakable}
\newtcolorbox{intuition}[1][Intuition]{
  colback=green!4, colframe=green!40!black, fonttitle=\bfseries,
  title=#1, breakable}
\newtcolorbox{warning}[1][Caution]{
  colback=red!4, colframe=red!50!black, fonttitle=\bfseries,
  title=#1, breakable}

% ---------- Code listings (Python by default) ----------
\lstdefinestyle{notes}{
  language=Python,
  basicstyle=\ttfamily\small,
  keywordstyle=\color{blue!70!black}\bfseries,
  commentstyle=\color{green!45!black}\itshape,
  stringstyle=\color{red!60!black},
  numbers=left, numberstyle=\tiny\color{gray},
  frame=single, rulecolor=\color{gray!40},
  backgroundcolor=\color{gray!5},
  breaklines=true, showstringspaces=false, tabsize=4,
  captionpos=b
}
\lstset{style=notes}

% ---------- Shared math notation ----------
% ============================================================
%  Shared math notation for ML / DL notes
% ============================================================

% ---------- Number sets ----------
\newcommand{\R}{\mathbb{R}}
\newcommand{\N}{\mathbb{N}}
\newcommand{\Z}{\mathbb{Z}}

% ---------- Probability ----------
\newcommand{\E}{\mathbb{E}}
\newcommand{\Prob}{\mathbb{P}}
\DeclareMathOperator{\Var}{Var}
\DeclareMathOperator{\Cov}{Cov}
\newcommand{\KL}[2]{D_{\mathrm{KL}}\!\left(#1 \,\|\, #2\right)}
\newcommand{\Normal}{\mathcal{N}}
\newcommand{\iid}{\overset{\text{i.i.d.}}{\sim}}

% ---------- Linear algebra ----------
\renewcommand{\vec}[1]{\bm{#1}}          % vectors: \vec{x}
\newcommand{\mat}[1]{\bm{#1}}            % matrices: \mat{W}
\newcommand{\T}{^{\top}}                 % transpose: \vec{x}\T
\DeclareMathOperator{\tr}{tr}
\DeclareMathOperator{\rank}{rank}
\DeclareMathOperator{\diag}{diag}
\DeclarePairedDelimiter{\norm}{\lVert}{\rVert}
\DeclarePairedDelimiter{\abs}{\lvert}{\rvert}
\DeclarePairedDelimiter{\inner}{\langle}{\rangle}

% ---------- Optimization / calculus ----------
\DeclareMathOperator*{\argmin}{arg\,min}
\DeclareMathOperator*{\argmax}{arg\,max}
\newcommand{\grad}{\nabla}
\newcommand{\pd}[2]{\frac{\partial #1}{\partial #2}}

% ---------- ML-specific ----------
\newcommand{\dataset}{\mathcal{D}}
\newcommand{\loss}{\mathcal{L}}
\newcommand{\hypo}{\mathcal{H}}
\newcommand{\params}{\bm{\theta}}
\newcommand{\X}{\mathcal{X}}
\newcommand{\Y}{\mathcal{Y}}
\newcommand{\indicator}[1]{\mathbb{1}\!\left[#1\right]}
\DeclareMathOperator{\sigmoid}{\sigma}
\DeclareMathOperator{\softmax}{softmax}
\DeclareMathOperator{\ReLU}{ReLU}
\DeclareMathOperator{\sign}{sign}

